%% sections/slide03-problem.tex — SLIDE 3 — Problem & objectives
%% =============================================================================
%% CONTENT BLUEPRINT — Problem statement + objectives
%% =============================================================================
%% Time budget: 45 seconds. This slide sets the stakes for the entire talk.
%%
%% WHAT TO PUT HERE:
%%
%%   Top half — Problem statement (exactly ONE paragraph, ≤ 2 sentences)
%%   ----------------------------------------------------------------------
%%   Format: "Despite [the gap shown in §1 of your thesis], [specific
%%   question X] remains [unquantified / unverified / unexplored] under
%%   [specific conditions Y]."
%%
%%   Examples (adjust to your project):
%%     • "Despite extensive work on Z-scan of bulk glasses, the wavelength
%%        dependence of n₂ in SF59 under nanosecond pulses has not been
%%        measured at room temperature for the 500–700 nm region."
%%     • "Although DFT-GGA sp-band calculations of TMD monolayers are
%%        abundant, no published value of the biaxial-strain coefficient of
%%        E_g is reproducible to within ±10 % for MoS₂ at 4 % strain."
%%
%%   DO NOT reproduce your §1.1 background here. Examiners have read your
%%   thesis abstract. This slide is the verdict, not the context.
%%
%%   Bottom half — Objectives (3–5 bullets)
%%   ---------------------------------------
%%   Each bullet must be: (a) ONE action verb, (b) testable, (c) achievable.
%%   Pattern: "To <verb> <object> <parameter range / sample set / method>."
%%
%%   Bad bullets:
%%     ✗ "To study non-linear optics."
%%     ✗ "To investigate materials."
%%     ✗ "To do simulations."
%%
%%   Good bullets:
%%     ✓ "To measure closed-aperture Z-scan traces of SF59 at λ = 632.8 nm
%%        and extract n₂ to ±15 %."
%%     ✓ "To compute the bandgap of monolayer MoS₂ under 0–8 % biaxial
%%        strain in steps of 1 % using PBE+SOC in Quantum ESPRESSO."
%%     ✓ "To compare computed absorption coefficients with published
%%        experimental values at hν = 1.5–3.0 eV."
%%
%%   Highlight the bullet that maps to the headline result of slide 5 in
%%   \highlight{...} (orange) so the examiner can trace Objective → Result.
%%
%% STYLE NOTES:
%%   • At most 2 short sentences for the problem statement.
%%   • 3–5 objectives, never more.
%%   • Never list deliverables (data DOI, code repo) on this slide — save
%%     those for slide 8 (summary).
%% =============================================================================

\begin{frame}{Statement of the Problem \& Objectives}
  \begin{beamercolorbox}[wd=.96\paperwidth, sep=8pt]{block body example}
    \textbf{\color{sustBlueDark}The problem.}%
    \medskip
    
    Despite extensive prior work on \highlight{[your topic]} in
    \highlight{[system class]}, no published value of \highlight{[target
    quantity]} \highlight{under [your conditions]} is reproducible to
    within $\pm$\,\highlight{[X]}\%.\hspace{1em}\emph{(Replace this paragraph
    with one sentence that states the specific gap your work fills.)}
  \end{beamercolorbox}

  \vspace{0.2cm}
  \textbf{\color{sustBlueDark}Objectives of this work.}
  \begin{itemize}\setlength\itemsep{2pt}
    \item \highlight{\textbf{O1}} \, To set up a [method] for measuring /
          computing [target quantity] on [system] at [conditions].
    \item \textbf{O2} \, To extract [intermediate observable] versus
          [parameter] at [range with N points].
    \item \textbf{O3} \, To compare measured / computed values against
          [X] published references covering [conditions].
    \item \textbf{O4} \, (Optional) To assess the role of [effect /
          approximation / correction].
  \end{itemize}

  \vspace{0.1cm}
  \timenote{Hit O1 motto in 5--7 s. Examiner will not remember every
  bullet; they need to remember the headline.}
\end{frame}
